%%%%%
%
%% Dans ce fichier, les �ventuels accents sont cod�s
%% en ISO-8859-15 (aka latin9)
%% La ligne suivante donne trois exemples 
%% e accent aigu : �, e accent grave : �, e-dans-l'o : �
\section{Calcul d'un flag qualit� pour la mesure de temp�rature dans les abris "fil fin" d�velopp�s par 4M}

\itemstyle{Op\'erateur~:} \operfmt{Calc\_flag\_filfin}
\index{Calc\_flag\_filfin}

%%%%%OBJET%%%%
\itemstyle{Objet~:}
Les abris dits "fil fin" pour la mesure de temp�rature ont �t� d�velopp�s par 4M pour la campagne Passy-2015.
L'�tude des donn�es de la campagne de test EFFIGI-2015 a montr� que l'incertitude de mesure de la temp�rature avec ces abris d�pend de l'orientation de l'ouverture de l'abri par rapport � la direction du vent et de la force du vent.
Le flag qualit� calcul� indique le niveau d'incertitude li� � l'abri :
\begin{center}
\begin{tabular}{|c|l|}
\hline
$Flag$ & Incertitude \\
\hline
1 & $0.15^\bullet C$ \\
2 & $0.30^\bullet C$\\
3 & $0.5^\bullet C$\\
\hline
\end{tabular}

\end{center}
%%%%%ARGUMENTS%%%%
\bigskip
%% Dans ce fichier, les �ventuels accents sont cod�s
%% en ISO-8859-15 (aka latin9)
%% La ligne suivante donne trois exemples 
%% e accent aigu : �, e accent grave : �, e-dans-l'o : �
%%%%%ENTREES%%%%
\begin{tabarg}
\multicolumn{3}{l}{\titretab {Entr\'ees}}\\
\hline
$FF$ & Vecteur, \texttt{float} & Force du vent \\
$DD$& Vecteur, \texttt{float} & Direction du vent (degr�s)\\
$Seuil_{FF_{AR}}$ & Coeff., \texttt{float} & Seuil de force du vent par vent arri�re (m\^eme unit� que $FF$)\\
$DD_{min}$ & Coeff., \texttt{float} & Borne inf�rieure du secteur vent de face (degr�s)\\
$DD_{max}$ & Coeff., \texttt{float} & Borne sup�rieure du secteur vent de face (degr�s)\\
\end{tabarg}

\medskip
%%%%%SORTIES%%%%
\begin{tabarg}
\multicolumn{3}{l}{\titretab {Sorties}}\\
\hline
$Flag$ & Vecteur, \texttt{float} & Flag qualit� de la temp�rature (unitless)\\
\end{tabarg}


%%%%%FORMULE%%%%
\itemstyle{Formule}

%% Formule de l'operateur : Calc_flag_filfin
%% Dans ce fichier, les �ventuels accents sont cod�s
%% en ISO-8859-15 (aka latin9)
%% La ligne suivante donne trois exemples 
%% e accent aigu : �, e accent grave : �, e-dans-l'o : �

\vspace{0.5cm}
$Flag =
\begin{cases}
1  \text{ si DD est dans l'intervalle }\left] DD_{min}, DD_{max} \right[ \text{ (cas du vent de face )
 ET  } FF \geq Seuil_{FF_{AR}} \\
2 \text{ si DD est dans l'intervalle }\left] DD_{min}, DD_{max} \right[ \text{ (cas du vent de face )
 ET  } FF < Seuil_{FF_{AR}} \\  
3 \text{ si DD est hors de l'intervalle } \left] DD_{min}, DD_{max} \right[ \text{ (cas du vent arri�re)}

\end{cases}
$





%%%%%AUTEUR%%%%
\itemstyle{Auteur~:}
CNRM/GMEI/TRAMM

%%%%%REFERENCE%%%%
\itemstyle{R\'ef\'erences~:}
Note de travail interne Sandrine Girres.
